% Generated draft: requires mathematical and source review.
\documentclass[11pt]{article}
\usepackage[margin=25mm]{geometry}
\usepackage{fontspec}
\setmainfont{Latin Modern Roman}
\usepackage{amsmath,amssymb,mathtools,unicode-math}
\setmathfont{Latin Modern Math}
\usepackage{enumitem,needspace,xurl,hyperref}
\setlength{\parindent}{0pt}
\setlength{\parskip}{5pt}
\setlength{\emergencystretch}{3em}
\newcommand{\N}{\mathbb{N}}
\newcommand{\Z}{\mathbb{Z}}
\newcommand{\Q}{\mathbb{Q}}
\newcommand{\R}{\mathbb{R}}
\newcommand{\C}{\mathbb{C}}
\newcommand{\F}{\mathbb{F}}
\newcommand{\A}{\mathbb{A}}
\newcommand{\PP}{\mathbb{P}}
\newcommand{\E}{\mathbb{E}}
\newcommand{\eps}{\varepsilon}
\newcommand{\dd}{\,\mathrm{d}}
\begin{document}
\section*{UnsolvedMath ID: 2305006}
\textbf{Source problem code:} catalogue entry
\subsection{Definitions and mathematical statement}
A function omits a set $A\subset\C$ if its image is disjoint from $A$. The source asks for conditions under which its formula labeled 5.9 remains valid when omitted values $w_n$ tend toward infinity sufficiently close to the imaginary axis. Neither formula 5.9 nor the ambient function class, normalization, domain, or coefficient notation $a_n$ is supplied. Thus the asserted conclusion cannot be written as an explicit inequality from this record.

The reported comparison cases are:
\begin{enumerate}
\item There exist functions omitting every value $2\pi i k$, $k\in\Z$, for which formula 5.9 fails. The source further reports $|a_n|\leq\log\log n$ for the relevant coefficients, citing Littlewood and Hayman, but supplies neither their definition nor a range of $n$. This displayed coefficient bound cannot be interpreted literally at all positive indices; no small-index or asymptotic convention is silently added.
\item Formula 5.9 holds if the function omits all of the imaginary axis outside some finite interval. In set notation, the omitted set includes $\{it:t\notin[a,b]\}$ for some finite real $a\leq b$. The source attributes this case to subordination.
\end{enumerate}
The proposed extension concerns omitted values clustering near infinity sufficiently close to the imaginary axis. The record does not specify a rate or inequality for $\operatorname{Re}w_n$, a distribution/separation condition, or the exact quantifier for ``sufficiently close.'' Absolute distance, relative distance and density along the axis give different hypotheses. No one of them is chosen, and no condition is invented to replace the missing conjecture.
\subsection{Short English statement}
Determine how strongly an omitted-value set must approach the imaginary axis at infinity to force the source's estimate 5.9. The record contrasts failure for the discrete omitted lattice $2\pi i\Z$ with validity when all but a bounded imaginary-axis interval is omitted. Because the estimate, function class and closeness criterion are absent, only that source-specific comparison can be stated. The supplied triage reports openness in Hayman's 2018 update and no recent resolution located; no independent status review is claimed.
\subsection{Sources}
\begin{itemize}
\item UnsolvedMath permanent ID 2305006, source code catalogue entry, ``Research Problems in Function Theory --- Problem 5.6.'' Registry curation and metadata: Ulam AI and UnsolvedMath Contributors, CC BY 4.0; updated 2026-08-06.
\item Hayman, \emph{Research Problems in Function Theory} (2018 source update), Problem 5.6, referring internally to formula 5.9: \url{https://arxiv.org/abs/1809.07200}. Supplied source-TeX extraction accessed 2026-07-29. The metadata retains the update for status triage and marks status/release and rights review NEEDS\_REVIEW.
\item Littlewood and Hayman are named for the omitted-lattice counterexamples and coefficient bound without complete bibliographic entries. No missing equation, title, date or locator is reconstructed by inference.
\end{itemize}
\end{document}
